\chapter{Symmetric Matrices and the Spectral Theorem}

In Chapter~24 you learned to take a matrix apart with its eigenvectors, and in
Chapter~25 you used that diagonalization to solve coupled differential
equations $\vec{u}' = \mat{A}\vec{u}$ using the matrix exponential
$e^{\mat{A}t} = \mat{P}e^{\mat{\Lambda}t}\mat{P}\inv$. The central idea was
decoupling: find special directions $\vec{x}$ where the matrix acts as pure
scaling ($\mat{A}\vec{x} = \lambda\vec{x}$), change to those eigenvector
coordinates so each mode evolves independently, and change back. But that
general story presented three complications. First, eigenvalues could be
complex, turning simple growth and decay into swirling spirals and
oscillations. Second, the eigenvectors could be skewed at awkward,
non-perpendicular angles, distorting geometry and forcing $\mat{P}\inv$ to be
an expensive matrix inverse. Third, defective matrices lacked enough
eigenvectors to diagonalize at all, spawning polynomial growth terms
$t e^{\lambda t}$ that tangled the modes.

This chapter is about the family of matrices where these
complications disappear: \textbf{symmetric matrices}. When a matrix
equals its own transpose -- when $\mat{A} = \mat{A}\T$ -- the situation changes
completely. Its eigenvalues are guaranteed real: no imaginary numbers, no
hidden rotations. Its eigenvectors can always be chosen \emph{perpendicular} to
one another, forming a rigid coordinate frame. And the change-of-basis matrix
$\mat{P}$ becomes an \textbf{orthogonal matrix} $\mat{Q}$ (a pure rotation or
reflection) whose inverse simplifies to $\mat{Q}\inv = \mat{Q}\T$.
The diagonalization collapses to
\[
  \boxed{\;\mat{A} = \mat{Q}\,\mat{\Lambda}\,\mat{Q}\T\;}
\]
The collection of eigenvalues of a matrix is called its \textbf{spectrum} --
named by analogy to the spectrum of light, where a glass prism breaks a
composite beam into its pure constituent colors. This factorization is the
\textbf{Spectral Theorem}. It decomposes a symmetric matrix into pure,
independent, perpendicular stretching actions. It underlies principal
component analysis in machine learning, the geometry of quadratic forms, the
vibrational modes of bridges and molecules, the covariance ellipsoids of
statistics, and the singular value decomposition. We start with the geometric
picture -- symmetric matrices stretch space along perpendicular axes -- and then
prove why symmetry forces both real eigenvalues and orthogonal eigenvectors.

\section{What Makes a Matrix Symmetric}
\label{ch26:what-is-symmetric}

A square matrix $\mat{A}$ is \textbf{symmetric} when it is unchanged by
transposition: $\mat{A}\T = \mat{A}$. Concretely, the entry in row $i$, column
$j$ equals the entry in row $j$, column $i$ -- the matrix is a mirror image of
itself across the main diagonal.

\[
  \mat{A} = \begin{bmatrix} 2 & 5 & -1 \\ 5 & 0 & 3 \\ -1 & 3 & 4 \end{bmatrix}
  \quad\text{is symmetric:}\quad a_{ij}=a_{ji}.
\]
Fold the matrix along the diagonal (the $2,0,4$) and the off-diagonal entries land
on their twins: the two $5$'s swap, the two $3$'s swap, the two $-1$'s swap. By
contrast
\[
  \mat{B} = \begin{bmatrix} 1 & 2 \\ -2 & 1 \end{bmatrix}
\]
is \emph{not} symmetric, because the $2$ above the diagonal faces a $-2$ below it.

\begin{definition}[Symmetric matrix]
\label{ch26:def-symmetric}
A real $n\times n$ matrix $\mat{A}$ is \textbf{symmetric} if $\mat{A}\T = \mat{A}$,
equivalently if $a_{ij} = a_{ji}$ for all $i,j$.
\end{definition}

\begin{focusbox}
Symmetry is a statement about the \emph{matrix}, but its real meaning is a
statement about the \emph{dot product}. The single identity to carry through this
whole chapter is
\[
  \inner{\mat{A}\vec{x}}{\vec{y}} = \inner{\vec{x}}{\mat{A}\vec{y}}
  \qquad\text{for all } \vec{x},\vec{y}.
\]
``You may move $\mat{A}$ from one side of the dot product to the other for free.''
Every theorem in this chapter falls out of that one sentence.
\end{focusbox}

Why is that the real meaning? Recall the dot product written with matrices:
$\inner{\vec{u}}{\vec{v}} = \vec{u}\T\vec{v}$. Then
\[
  \inner{\mat{A}\vec{x}}{\vec{y}}
  = (\mat{A}\vec{x})\T\vec{y}
  = \vec{x}\T \mat{A}\T \vec{y}
  = \vec{x}\T(\mat{A}\T\vec{y})
  = \inner{\vec{x}}{\mat{A}\T\vec{y}}.
\]
So $\mat{A}$ can always be moved across the dot product if you transpose it on the
way. Symmetry, $\mat{A}\T=\mat{A}$, is exactly the condition that it crosses
\emph{without changing}. Hold onto this; we will use it four times.

\begin{intuitionbox}
A general matrix can rotate, stretch, and shear space. A symmetric matrix is
much simpler: geometrically it does nothing but \textbf{stretch space along a set
of perpendicular axes}. It introduces no twist or non-orthogonal shear between the
principal axes. This action corresponds to stretching an elastic membrane along
perpendicular axes with different scaling factors.
\end{intuitionbox}

\subsection{Where symmetric matrices come from}
\label{ch26:where-from}

You might worry that symmetric matrices are a rare, hand-picked species. The
opposite is true -- they are everywhere, and usually for the same structural
reason: a symmetric matrix records a \emph{relationship between pairs} that does
not care about order.

\begin{itemize}
  \item \textbf{Gram matrices.} For \emph{any} matrix $\mat{M}$, the product
        $\mat{M}\T\mat{M}$ is symmetric, because
        $(\mat{M}\T\mat{M})\T = \mat{M}\T(\mat{M}\T)\T = \mat{M}\T\mat{M}$.
        Its $(i,j)$ entry is the dot product of column $i$ with column $j$ of
        $\mat{M}$ -- and the dot product does not care about order. These appear
        in least squares as the normal equations $\mat{M}\T\mat{M}\vec{x} =
        \mat{M}\T\vec{b}$.
  \item \textbf{Covariance matrices.} A covariance matrix measures how pairs of
        random variables vary together. Because the co-movement of feature $i$
        with feature $j$ is identical to that of $j$ with $i$, every covariance
        matrix in statistics and machine learning is symmetric. (As you will
        see, this is simply the Gram matrix of the mean-centered data, up to a scalar factor.)
  \item \textbf{Distances and similarities.} The matrix of pairwise distances or
        similarities between $n$ objects is symmetric: object $i$ is as far from
        $j$ as $j$ is from $i$.
  \item \textbf{Second derivatives.} The \textbf{Hessian matrix} of second
        partial derivatives, $H_{ij} = \partial^2 f / \partial x_i \partial x_j$,
        captures the local curvature of a multivariable surface in every
        direction. It is symmetric because mixed partial derivatives commute
        (Clairaut's theorem: $\frac{\partial^2 f}{\partial x_i \partial x_j} =
        \frac{\partial^2 f}{\partial x_j \partial x_i}$). Modern optimization
        algorithms rely on the Hessian to curve toward minima.
  \item \textbf{Graphs and networks.} The adjacency matrix of an undirected graph
        is symmetric: an edge between node $i$ and node $j$ is equally an edge
        between $j$ and $i$. So is the \textbf{graph Laplacian} $\mat{L} =
        \mat{D} - \mat{A}$ (the diagonal matrix of node connection counts minus
        the adjacency matrix), which governs how heat, random walks, and
        information diffuse across a network.
\end{itemize}

Whenever the underlying relationship is undirected -- a pairing with no preferred
arrow from one party to the other -- the matrix that records it is symmetric.
That is why this chapter pays off so widely.

\section{The Picture: Stretching Along Perpendicular Axes}
\label{ch26:picture}

Before any theorem, let us look. Take the symmetric matrix
\[
  \mat{A} = \begin{bmatrix} 2 & 1 \\ 1 & 2 \end{bmatrix},
\]
and watch what it does to the unit circle. The circle is the set of all arrows of
length $1$; its image under $\mat{A}$ tells you everything about how $\mat{A}$
stretches space. Figure~\ref{ch26:fig-ellipse} shows the result: the circle
becomes an \textbf{ellipse}, and -- this is the whole point -- the long and short
axes of that ellipse are \emph{perpendicular}, and they line up with the diagonal
directions $(1,1)$ and $(1,-1)$.

\begin{figure}[H]
\centering
\begin{tikzpicture}[scale=1.1,>=Stealth]
  % left: unit circle with the two special directions
  \begin{scope}[shift={(0,0)}]
    \draw[->,gray!55] (-1.7,0)--(1.9,0) node[right]{$x$};
    \draw[->,gray!55] (0,-1.7)--(0,1.9) node[above]{$y$};
    \draw[vecblue,thick] (0,0) circle (1);
    \draw[->,vecred,thick] (0,0)--(0.7071,0.7071) node[above right]{$\vec{q}_1$};
    \draw[->,structurepurple,thick] (0,0)--(0.7071,-0.7071) node[below right]{$\vec{q}_2$};
    \node at (0,-2.85) {\small unit circle};
  \end{scope}
  % arrow
  \node at (2.7,0) {\Large $\xrightarrow{\;\mat{A}\;}$};
  % right: image ellipse, axes scaled by 3 and 1
  \begin{scope}[shift={(5.6,0)}]
    \draw[->,gray!55] (-2.4,0)--(3.4,0) node[right]{$x$};
    \draw[->,gray!55] (0,-2.4)--(0,2.6) node[above]{$y$};
    % ellipse with semi-axes 3 (along (1,1)) and 1 (along (1,-1)), rotated 45deg
    \draw[targetgreen,thick,rotate=45] (0,0) ellipse (3 and 1);
    \draw[->,vecred,thick] (0,0)--(2.1213,2.1213) node[above right]{$3\vec{q}_1$};
    \draw[->,structurepurple,thick] (0,0)--(0.7071,-0.7071) node[below right=2pt,fill=white,inner sep=1pt]{$1\vec{q}_2$};
    \node at (0,-2.85) {\small image: an ellipse};
  \end{scope}
\end{tikzpicture}
\caption{The symmetric matrix $\mat{A}=\protect\begin{bmatrix}2&1\\1&2\protect\end{bmatrix}$
sends the unit circle to an ellipse. The special directions
$\vec{q}_1 \propto (1,1)$ and $\vec{q}_2 \propto (1,-1)$ are
\emph{perpendicular}, and $\mat{A}$ simply stretches them by $\lambda_1 = 3$ and
$\lambda_2 = 1$. These are the \textbf{principal axes}.}
\label{ch26:fig-ellipse}
\end{figure}

Let us verify the picture by hand. Along $(1,1)$:
\[
  \mat{A}\colvec{1\\1} = \colvec{2+1\\1+2} = \colvec{3\\3} = 3\colvec{1\\1}.
\]
So $(1,1)$ is an eigenvector with eigenvalue $3$ -- $\mat{A}$ stretches it by a
factor of $3$. Along $(1,-1)$:
\[
  \mat{A}\colvec{1\\-1} = \colvec{2-1\\1-2} = \colvec{1\\-1} = 1\cdot\colvec{1\\-1}.
\]
So $(1,-1)$ is an eigenvector with eigenvalue $1$ -- $\mat{A}$ leaves its length
alone. The two eigenvectors are perpendicular: $(1,1)\cdot(1,-1) = 1-1 = 0$.

That perpendicularity is not luck. It is forced by symmetry, and it is the heart
of the spectral theorem. The long and short axes of the ellipse -- the directions
along which a symmetric matrix purely stretches -- are called the \textbf{principal
axes}, and they are always at right angles.

\begin{walkbox}
Consider the matrix $\mat{A}=\begin{bmatrix}2&1\\1&2\end{bmatrix}$. We
found it stretches $(1,1)$ by $3$ and fixes $(1,-1)$. Now picture feeding it a
non-eigenvector like $(1,0)$. It must land between the scaling along the two axes.
Which way does the output arrow tilt -- toward the dominant axis or the smaller one?
The output always leans toward the direction of \emph{larger} eigenvalue.
Checking directly: $\mat{A}(1,0)=(2,1)$, which tilts toward the $(1,1)$ axis.
\end{walkbox}

\section{Real Eigenvalues}
\label{ch26:real-eigenvalues}

A general matrix can have complex eigenvalues -- a rotation matrix has eigenvalues
$\cos\theta \pm i\sin\theta$, which live off the real line and correspond to no
real stretching direction at all. Symmetric matrices never do this. Their
eigenvalues are always real numbers, every time.

\begin{theorem}[Symmetric matrices have real eigenvalues]
\label{ch26:thm-real}
If $\mat{A}$ is a real symmetric matrix, then every eigenvalue of $\mat{A}$ is
real.
\end{theorem}

The cleanest proof briefly steps into complex numbers, uses the conjugate, and
steps back out. It is three lines once you see the trick.
We use the \textbf{conjugate transpose} $\vec{z}^{*} = \overline{\vec{z}}\T$:
swap rows and columns while replacing every imaginary number $i$ with $-i$. In
ordinary real geometry, $\vec{x}\T\vec{x} = \sum x_k^2$ computes the squared
length of a vector; for complex vectors, the conjugate ensures that
$\vec{z}^{*}\vec{z} = \sum_k |z_k|^2$ is a \emph{positive real number} whenever
$\vec{z}\neq\vec{0}$.

\begin{proof}
Let $\lambda$ be an eigenvalue (a priori possibly complex) with eigenvector
$\vec{z}\neq\vec{0}$, so $\mat{A}\vec{z} = \lambda\vec{z}$. Form the scalar
$\vec{z}^{*}\mat{A}\vec{z}$ and compute it two ways.

\emph{First way.} Substitute $\mat{A}\vec{z}=\lambda\vec{z}$:
\[
  \vec{z}^{*}\mat{A}\vec{z} = \vec{z}^{*}(\lambda\vec{z}) = \lambda\,(\vec{z}^{*}\vec{z}).
\]

\emph{Second way.} Take the conjugate transpose of the whole scalar. By conjugate transpose rules,
$(\vec{z}^{*}\mat{A}\vec{z})^{*} = \vec{z}^{*}\mat{A}^{*}\vec{z}$. Because $\mat{A}$
is \emph{real} and \emph{symmetric}, $\mat{A}^{*} = \overline{\mat{A}}\T = \mat{A}\T = \mat{A}$,
so $(\vec{z}^{*}\mat{A}\vec{z})^{*} = \vec{z}^{*}\mat{A}\vec{z}$. Therefore
\[
  \vec{z}^{*}\mat{A}\vec{z}
  = (\mat{A}\vec{z})^{*}\vec{z}
  = (\lambda\vec{z})^{*}\vec{z}
  = \overline{\lambda}\,(\vec{z}^{*}\vec{z}).
\]

Setting the two equal: $\lambda\,(\vec{z}^{*}\vec{z}) = \overline{\lambda}\,(\vec{z}^{*}\vec{z})$.
Since $\vec{z}^{*}\vec{z} > 0$, we may cancel it to get $\lambda = \overline{\lambda}$.
A number equal to its own conjugate is real.
\end{proof}

\begin{intuitionbox}
A complex eigenvalue corresponds to rotation combined with scaling. A symmetric
matrix satisfies $\inner{\mat{A}\vec{x}}{\vec{y}} = \inner{\vec{x}}{\mat{A}\vec{y}}$,
which prevents rotational coupling across orthogonal directions and restricts the
transformation to pure directional scaling. Real eigenvalues are the algebraic
signature of this purely scaling behavior.
\end{intuitionbox}

\section{Orthogonal Eigenvectors}
\label{ch26:orthogonal-eigenvectors}

The second fundamental property of symmetric matrices is that
eigenvectors belonging to \emph{different} eigenvalues are automatically
perpendicular. The proof relies directly on the symmetry identity
$\inner{\mat{A}\vec{x}}{\vec{y}} = \inner{\vec{x}}{\mat{A}\vec{y}}$.

\begin{theorem}[Eigenvectors for distinct eigenvalues are orthogonal]
\label{ch26:thm-orthogonal}
Let $\mat{A}$ be real symmetric. If $\mat{A}\vec{x}=\lambda\vec{x}$ and
$\mat{A}\vec{y}=\mu\vec{y}$ with $\lambda\neq\mu$, then $\inner{\vec{x}}{\vec{y}}=0$.
\end{theorem}

\begin{proof}
Compute the single number $\inner{\mat{A}\vec{x}}{\vec{y}}$ two ways, sliding
$\mat{A}$ across the dot product using symmetry.

On one hand, $\mat{A}\vec{x}=\lambda\vec{x}$, so
\[
  \inner{\mat{A}\vec{x}}{\vec{y}} = \inner{\lambda\vec{x}}{\vec{y}} = \lambda\inner{\vec{x}}{\vec{y}}.
\]
On the other hand, by symmetry $\inner{\mat{A}\vec{x}}{\vec{y}}=\inner{\vec{x}}{\mat{A}\vec{y}}$,
and $\mat{A}\vec{y}=\mu\vec{y}$, so
\[
  \inner{\mat{A}\vec{x}}{\vec{y}} = \inner{\vec{x}}{\mat{A}\vec{y}} = \inner{\vec{x}}{\mu\vec{y}} = \mu\inner{\vec{x}}{\vec{y}}.
\]
Subtracting, $(\lambda-\mu)\inner{\vec{x}}{\vec{y}} = 0$. Since $\lambda\neq\mu$,
the factor $\lambda-\mu\neq 0$, so $\inner{\vec{x}}{\vec{y}}=0$.
\end{proof}

That is the entire argument: push $\mat{A}$ onto the eigenvector on the right, pull
it onto the eigenvector on the left, and the two eigenvalues must disagree unless
the vectors are perpendicular. Three lines, and the perpendicular axes of every
ellipse in the chapter are explained.

\begin{example}[Checking orthogonality numerically]
\label{ch26:ex-ortho}
Return to $\mat{A}=\begin{bmatrix}2&1\\1&2\end{bmatrix}$ with eigenvectors
$\vec{x}=(1,1)$ for $\lambda=3$ and $\vec{y}=(1,-1)$ for $\mu=1$. Directly checking
the dot product: $\inner{\vec{x}}{\vec{y}} = (1)(1)+(1)(-1) = 0$. Now a $3\times 3$ example: take
\[
  \mat{B} = \begin{bmatrix} 1 & 2 & 0 \\ 2 & 1 & 0 \\ 0 & 0 & 5 \end{bmatrix}.
\]
The bottom-right block already exposes the eigenvector $\vec{e}_3=(0,0,1)$ with
$\lambda=5$. The top-left $2\times 2$ block is $\begin{bmatrix}1&2\\2&1\end{bmatrix}$,
whose eigenvectors are $(1,1,0)$ (eigenvalue $3$) and $(1,-1,0)$ (eigenvalue
$-1$). All three pairwise dot products vanish:
\[
  (1,1,0)\cdot(1,-1,0)=0,\quad (1,1,0)\cdot(0,0,1)=0,\quad (1,-1,0)\cdot(0,0,1)=0.
\]
Three mutually perpendicular stretching axes, exactly as promised.
\end{example}

\begin{recoverybox}
\textbf{Pause and think: the third direction for free.} Suppose a $3\times 3$
symmetric matrix has three distinct eigenvalues, and you have already computed
eigenvectors $\vec{x}_1 = (1, 1, 0)$ and $\vec{x}_2 = (1, -1, 0)$ for the
first two. Do you need to set up and solve a third system of equations
$(\mat{A} - \lambda_3\mat{I})\vec{x} = \vec{0}$ to find the direction of the
third eigenvector?

No! Because $\mat{A}$ is symmetric and the eigenvalues are distinct,
Theorem~\ref{ch26:thm-orthogonal} guarantees that the third eigenvector
$\vec{x}_3$ must be perpendicular to both $\vec{x}_1$ and $\vec{x}_2$. In
$\R^3$, the only vectors perpendicular to the entire $xy$-plane point along the
$z$-axis, so $\vec{x}_3$ must be proportional to $(0, 0, 1)$ (or their cross
product $\vec{x}_1 \times \vec{x}_2$). Orthogonality hands you the final axis
with zero row reduction.
\end{recoverybox}

\begin{trapbox}
Theorem~\ref{ch26:thm-orthogonal} only guarantees orthogonality across
\emph{distinct} eigenvalues. If an eigenvalue is repeated -- say a double
eigenvalue with a whole \emph{plane} of eigenvectors -- two eigenvectors from that
same plane need not start out perpendicular. The fix is easy and important: within
a repeated eigenvalue's eigenspace you are free to \emph{choose} an orthonormal
basis (Gram--Schmidt does it). So you can \emph{always} end up with a full
orthonormal set; you just may have to pick wisely inside a shared eigenspace. The
full spectral theorem in the next section asserts this can always be done.
\end{trapbox}

\section{The Spectral Theorem}
\label{ch26:spectral-theorem}

Combining real eigenvalues (Theorem~\ref{ch26:thm-real}) with an
orthonormal basis of eigenvectors yields the fundamental factorization
of symmetric matrices.

\begin{formalbox}
\begin{theorem}[Spectral Theorem for real symmetric matrices]
\label{ch26:thm-spectral}
Let $\mat{A}$ be a real $n\times n$ symmetric matrix. Then there is an
\textbf{orthogonal} matrix $\mat{Q}$ (its columns form an orthonormal basis of
$\R^n$, so $\mat{Q}\T\mat{Q}=\mat{I}$ and $\mat{Q}\inv=\mat{Q}\T$) and a
\textbf{real diagonal} matrix $\mat{\Lambda}=\diag(\lambda_1,\dots,\lambda_n)$
such that
\[
  \mat{A} = \mat{Q}\,\mat{\Lambda}\,\mat{Q}\T.
\]
The columns $\vec{q}_1,\dots,\vec{q}_n$ of $\mat{Q}$ are orthonormal eigenvectors
of $\mat{A}$, and $\lambda_1,\dots,\lambda_n$ are the corresponding (real)
eigenvalues. Equivalently, every real symmetric matrix is \textbf{orthogonally
diagonalizable}.
\end{theorem}
\end{formalbox}

The single most important structural change from ordinary diagonalization
$\mat{A}=\mat{P}\mat{\Lambda}\mat{P}\inv$ is that $\mat{P}\inv$ has become
$\mat{Q}\T$. The hard, expensive inverse turns into a free transpose, because the
eigenvector basis is orthonormal. That one swap is what makes the spectral
theorem so computationally and conceptually pleasant.

\subsection{Reading the factorization as a sum of projections}
\label{ch26:sum-of-projections}

Multiply $\mat{Q}\mat{\Lambda}\mat{Q}\T$ out column-by-column and a cleaner form
appears. Writing $\mat{Q}=[\vec{q}_1\ \cdots\ \vec{q}_n]$, the product expands into
\[
  \mat{A} = \mat{Q}\,\mat{\Lambda}\,\mat{Q}\T
  = \sum_{k=1}^{n} \lambda_k\,\vec{q}_k\vec{q}_k\T.
\]
Each term $\vec{q}_k\vec{q}_k\T$ is an $n\times n$ matrix of rank one -- it is the
\textbf{projection} onto the line through $\vec{q}_k$ (recall the projection onto a
unit vector $\vec{q}$ is $\vec{q}\vec{q}\T$). So the spectral theorem reads:

\begin{focusbox}
\[
  \mat{A} = \lambda_1\,\vec{q}_1\vec{q}_1\T + \lambda_2\,\vec{q}_2\vec{q}_2\T + \cdots + \lambda_n\,\vec{q}_n\vec{q}_n\T.
\]
A symmetric matrix is a \textbf{weighted sum of perpendicular projections}. To
apply $\mat{A}$ to a vector: project onto each principal axis, scale that piece by
the axis's eigenvalue, and add the pieces back up. That is \emph{all} a symmetric
matrix does.
\end{focusbox}

Figure~\ref{ch26:fig-projsum} makes this concrete for
$\mat{A}=\begin{bmatrix}2&1\\1&2\end{bmatrix}$ acting on $\vec{x}=(1,0)$: split
$\vec{x}$ into its components along the two perpendicular axes, scale each by its
eigenvalue, and add the scaled pieces to land exactly on $\mat{A}\vec{x}=(2,1)$.

\begin{figure}[H]
\centering
\begin{tikzpicture}[scale=1.5,>=Stealth]
  \draw[->,gray!55] (-0.4,0)--(2.4,0) node[right]{$x$};
  \draw[->,gray!55] (0,-1.0)--(0,1.9) node[above]{$y$};
  % faint principal-axis directions
  \draw[gray!40,dashed] (-0.5,-0.5)--(2.1,2.1) node[above right,gray!60]{\small$\vec{q}_1$};
  \draw[gray!40,dashed] (-0.7,0.7)--(1.1,-1.1) node[below right,gray!60]{\small$\vec{q}_2$};
  % input x=(1,0)
  \draw[->,vecblue,very thick] (0,0)--(1,0) node[below right]{\small$\vec{x}=(1,0)$};
  % components p1=(0.5,0.5), p2=(0.5,-0.5)
  \draw[->,vecred,thick] (0,0)--(0.5,0.5) node[above left]{\small$\vec{p}_1$};
  \draw[->,structurepurple,thick] (0,0)--(0.5,-0.5) node[below left]{\small$\vec{p}_2$};
  % scaled: 3 p1 = (1.5,1.5), 1 p2 = (0.5,-0.5); chain them to Ax=(2,1)
  \draw[->,vecred,thick,dotted] (0,0)--(1.5,1.5) node[above left]{\small$\lambda_1\vec{p}_1$};
  \draw[->,structurepurple,thick,dotted] (1.5,1.5)--(2.0,1.0) node[above right,midway]{\small$+\lambda_2\vec{p}_2$};
  \draw[->,targetgreen,very thick] (0,0)--(2.0,1.0) node[right]{\small$\mat{A}\vec{x}=(2,1)$};
\end{tikzpicture}
\caption{The spectral decomposition in action on
$\mat{A}=\protect\begin{bmatrix}2&1\\1&2\protect\end{bmatrix}$, $\vec{x}=(1,0)$.
First split $\vec{x}$ into its perpendicular components
$\vec{p}_1=\vec{q}_1\vec{q}_1\T\vec{x}=(\tfrac12,\tfrac12)$ along the strong axis
and $\vec{p}_2=\vec{q}_2\vec{q}_2\T\vec{x}=(\tfrac12,-\tfrac12)$ along the weak
axis (so $\vec{p}_1+\vec{p}_2=\vec{x}$). Then scale by the eigenvalues
$\lambda_1=3$ and $\lambda_2=1$ and add: $3\vec{p}_1+1\vec{p}_2=(2,1)=\mat{A}\vec{x}$.
``Project, scale, add'' is the whole machine.}
\label{ch26:fig-projsum}
\end{figure}

\begin{recoverybox}
\textbf{Pause and think: testing the projection sum.} What happens if you
multiply the sum $\sum_{k=1}^n \lambda_k \vec{q}_k\vec{q}_k\T$ directly by
the first eigenvector $\vec{q}_1$?

Because the basis is orthonormal, $\vec{q}_1\T\vec{q}_1 = 1$ while
$\vec{q}_k\T\vec{q}_1 = 0$ for every $k \neq 1$. Every term in the sum vanishes
except the first, leaving $\lambda_1\vec{q}_1(1) = \lambda_1\vec{q}_1$. The
eigenvector equation is hard-wired directly into the sum of projections.
\end{recoverybox}

This decomposition is called the \textbf{spectral decomposition}, and it is the
direct ancestor of the singular value decomposition. When you later drop the
smallest $|\lambda_k|$ terms to compress a matrix, you are doing the linear algebra
behind PCA and low-rank approximation.

\subsection{The geometry: rotate, stretch, rotate back}
\label{ch26:rotate-stretch}

Read $\mat{A}\vec{x} = \mat{Q}\mat{\Lambda}\mat{Q}\T\vec{x}$ right-to-left as three
moves, just as you read ordinary diagonalization -- but now the first and last
moves are rigid.

\begin{enumerate}
  \item $\mat{Q}\T\vec{x}$: \textbf{rotate} (or reflect) space so the principal
        axes line up with the coordinate axes. This step measures how much of
        $\vec{x}$ points along each $\vec{q}_k$; its $k$-th entry is
        $\inner{\vec{q}_k}{\vec{x}}$. Because $\mat{Q}$ is orthogonal, lengths and
        angles are perfectly preserved -- nothing is distorted yet.
  \item $\mat{\Lambda}(\cdots)$: \textbf{stretch} each coordinate axis by its
        eigenvalue $\lambda_k$ (a negative $\lambda_k$ flips that axis).
  \item $\mat{Q}(\cdots)$: \textbf{rotate back} to the original orientation.
\end{enumerate}

\begin{figure}[H]
\centering
\begin{tikzpicture}[scale=0.62,>=Stealth]
  % stage 1: tilted axes with a unit circle
  \begin{scope}[shift={(0,0)}]
    \draw[gray!35] (-1.5,-1.5) grid[step=0.5] (1.5,1.5);
    \draw[vecblue,thick] (0,0) circle (1);
    \draw[->,vecred,thick,rotate=30] (0,0)--(1.0,0);
    \draw[->,structurepurple,thick,rotate=30] (0,0)--(0,1.0);
    \node at (0,-2.1) {\small original axes $\vec{q}_k$};
  \end{scope}
  \node at (2.6,0) {\small$\xrightarrow{\;\mat{Q}\T\;}$};
  % stage 2: aligned axes, still a circle
  \begin{scope}[shift={(5.2,0)}]
    \draw[gray!35] (-1.5,-1.5) grid[step=0.5] (1.5,1.5);
    \draw[vecblue,thick] (0,0) circle (1);
    \draw[->,vecred,thick] (0,0)--(1.0,0);
    \draw[->,structurepurple,thick] (0,0)--(0,1.0);
    \node at (0,-2.1) {\small axes aligned};
  \end{scope}
  \node at (7.8,0) {\small$\xrightarrow{\;\mat{\Lambda}\;}$};
  % stage 3: stretched ellipse, axis-aligned
  \begin{scope}[shift={(10.8,0)}]
    \draw[gray!35] (-1.5,-1.5) grid[step=0.5] (2.0,1.5);
    \draw[targetgreen,thick] (0,0) ellipse (1.8 and 0.7);
    \draw[->,vecred,thick] (0,0)--(1.8,0);
    \draw[->,structurepurple,thick] (0,0)--(0,0.7);
    \node at (0.3,-2.1) {\small stretched};
  \end{scope}
  \node at (13.6,0) {\small$\xrightarrow{\;\mat{Q}\;}$};
  % stage 4: rotated back
  \begin{scope}[shift={(16.8,0)}]
    \draw[gray!35] (-1.5,-1.5) grid[step=0.5] (2.0,1.5);
    \draw[targetgreen,thick,rotate=30] (0,0) ellipse (1.8 and 0.7);
    \draw[->,vecred,thick,rotate=30] (0,0)--(1.8,0);
    \draw[->,structurepurple,thick,rotate=30] (0,0)--(0,0.7);
    \node at (0.3,-2.1) {\small rotated back};
  \end{scope}
\end{tikzpicture}
\caption{The action $\mat{A}=\mat{Q}\mat{\Lambda}\mat{Q}\T$ read right to left.
$\mat{Q}\T$ rigidly rotates the principal axes onto the coordinate axes (lengths
unchanged), $\mat{\Lambda}$ stretches each coordinate axis by its eigenvalue, and a
final $\mat{Q}$ rotates back. A symmetric matrix is ``rotate, stretch, rotate
back'' -- with the \emph{same} rotation undone at the end.}
\label{ch26:fig-rotate-stretch}
\end{figure}

\begin{strangbox}
Strang calls this $\mat{A} = \mat{Q}\mat{\Lambda}\mat{Q}\T$ and sometimes
``$\mat{S} = \mat{Q}\mat{\Lambda}\mat{Q}\T$,'' reserving $\mat{S}$ for symmetric
matrices. He summarizes it in three words: \emph{real eigenvalues, orthogonal
eigenvectors}. \textbf{Before} Strang, the payoff to hold onto is the picture in
Figure~\ref{ch26:fig-rotate-stretch}: symmetric means perpendicular stretching
axes. \textbf{After} Strang, notice that the SVD ($\mat{A}=\mat{U}\mat{\Sigma}\mat{V}\T$)
is the same idea generalized to \emph{rectangular} matrices, where the input
rotation $\mat{V}$ and the output rotation $\mat{U}$ are allowed to differ. When $\mat{A}$ is positive semidefinite, the spectral theorem is the SVD's symmetric special case, with $\mat{U}=\mat{V}=\mat{Q}$ and $\mat{\Sigma}=\mat{\Lambda}$ (more generally $\mat{\Sigma}=\abs{\mat{\Lambda}}$, with $\mat{U}$ and $\mat{V}$ differing by sign flips on directions with negative eigenvalues).
\end{strangbox}

\subsection{Why it is true: a proof sketch}
\label{ch26:proof-sketch}

We have already proved the two pillars: eigenvalues are real
(Theorem~\ref{ch26:thm-real}) and eigenvectors for distinct eigenvalues are
orthogonal (Theorem~\ref{ch26:thm-orthogonal}). The only gap is \emph{repeated}
eigenvalues -- we must be sure there are always enough orthonormal eigenvectors to
fill an entire orthonormal basis, even when an eigenvalue repeats. The argument
runs by induction on the matrix size $n$, using a technique known as
\textbf{deflation}: peel off one known eigenvector $\vec{q}_1$, and show that
the remaining problem drops cleanly into the perpendicular subspace of one
smaller dimension, without ever leaking back.

\begin{proof}[Proof sketch]
The case $n=1$ is trivial ($\mat{A}$ is a single real number). Suppose the theorem
holds for all symmetric matrices of size $n-1$. Given a symmetric $n\times n$
matrix $\mat{A}$, a real eigenvalue $\lambda_1$ exists (the characteristic
polynomial has a real root, by Theorem~\ref{ch26:thm-real}), with a \emph{unit}
eigenvector $\vec{q}_1$: $\mat{A}\vec{q}_1=\lambda_1\vec{q}_1$.

Now look at the \textbf{orthogonal complement} $W = \vec{q}_1^{\perp}$, the
$(n-1)$-dimensional subspace of all vectors perpendicular to $\vec{q}_1$. The key
claim is that $\mat{A}$ \emph{maps $W$ into itself}. To check it, take any
$\vec{w}\in W$, so $\inner{\vec{q}_1}{\vec{w}}=0$, and test whether $\mat{A}\vec{w}$
is still perpendicular to $\vec{q}_1$:
\[
  \inner{\vec{q}_1}{\mat{A}\vec{w}}
  = \inner{\mat{A}\vec{q}_1}{\vec{w}}
  = \inner{\lambda_1\vec{q}_1}{\vec{w}}
  = \lambda_1\inner{\vec{q}_1}{\vec{w}} = 0.
\]
(We slid $\mat{A}$ across the dot product using symmetry -- the same move again.)
So $\mat{A}\vec{w}\in W$. Restricting $\mat{A}$ to the $(n-1)$-dimensional space
$W$ gives a smaller symmetric operator, which by the induction hypothesis has an
orthonormal eigenbasis $\vec{q}_2,\dots,\vec{q}_n$ of $W$. Together with
$\vec{q}_1$, these form an orthonormal eigenbasis of all of $\R^n$. Stacking them
as columns of $\mat{Q}$ gives $\mat{A}\mat{Q}=\mat{Q}\mat{\Lambda}$, i.e.
$\mat{A}=\mat{Q}\mat{\Lambda}\mat{Q}\T$.
\end{proof}

\begin{figure}[H]
\centering
\begin{tikzpicture}[scale=1.1,>=Stealth]
  % Shaded plane W = q1^\perp drawn with 3D perspective
  \fill[sourceblue!10,draw=sourceblue!70,thick] (-2.8,-1.1) -- (1.4,-1.1) -- (2.8,1.1) -- (-1.4,1.1) -- cycle;
  \node[sourceblue!90!black] at (-1.1,0.7) {\small $W = (\vec{q}_1)^\perp$};

  % origin
  \coordinate (O) at (0,0);

  % right-angle marker between q1 and W
  \draw[gray!60] (0,0.3) -- (0.3,0.3) -- (0.3,0);

  % vector w in W
  \draw[->,quotientorange,very thick] (O) -- (1.7,-0.5) node[below right]{\small $\vec{w}$};
  % A w stays in W
  \draw[->,targetgreen,very thick] (O) -- (-1.3,-0.6) node[below left]{\small $\mat{A}\vec{w}\in W$};

  % eigenvectors q2, q3 inside W found by induction
  \draw[->,structurepurple,thick] (O) -- (1.9,0.5) node[above right]{\small $\vec{q}_2$};
  \draw[->,structurepurple,thick] (O) -- (-1.6,0.35) node[below left]{\small $\vec{q}_3$};

  % vector q1 perpendicular to W along z-axis
  \draw[->,vecred,very thick] (O) -- (0,1.8) node[left]{\small $\vec{q}_1$};
  \draw[->,vecred,dashed,thick] (0,1.8) -- (0,2.6) node[above]{\small $\mat{A}\vec{q}_1 = \lambda_1\vec{q}_1$};
\end{tikzpicture}
\caption{The geometric engine of the Spectral Theorem: deflation. Once the first unit eigenvector $\vec{q}_1$ is found, symmetry forces $\mat{A}$ to map the perpendicular slice $W = \vec{q}_1^\perp$ strictly back into itself ($\inner{\vec{q}_1}{\mat{A}\vec{w}} = \lambda_1\inner{\vec{q}_1}{\vec{w}} = 0$). The transformation never ``leaks'' out of $W$. We can therefore restrict $\mat{A}$ to $W$ and repeat the argument by induction, finding an orthonormal eigenbasis $\vec{q}_2,\dots,\vec{q}_n$ within the smaller subspace.}
\label{ch26:fig-deflation}
\end{figure}

The engine of the whole proof is ``peel off one eigenvector, and symmetry
guarantees the rest of the matrix still lives cleanly in the perpendicular
subspace'' (Figure~\ref{ch26:fig-deflation}). Symmetry is exactly what makes the
orthogonal complement \emph{invariant}, and that is what lets the induction go
through even with repeated eigenvalues.

\begin{aestheticbox}
The spectral theorem is among the most consequential theorems in applied mathematics.
It shows that coupled pairwise relationships -- covariances, quadratic energies,
distances, and physical couplings -- decompose into a set of independent,
mutually perpendicular one-dimensional actions. Diagonalization decouples the system
and reveals the natural axes along which the transformation acts by pure scaling.
Vibrating structures possess normal modes; data clouds align along principal components;
quantum observables admit stationary energy states. Each is an instance of the spectral
decomposition applied in a different mathematical or physical setting.
\end{aestheticbox}

\section{Computing the Spectral Decomposition by Hand}
\label{ch26:computing}

Let us run the full procedure on a concrete matrix so the abstract $\mat{Q}$,
$\mat{\Lambda}$ become numbers you can hold.

\begin{example}[Full spectral decomposition of a $2\times2$ matrix]
\label{ch26:ex-full-2x2}
Take
\[
  \mat{A} = \begin{bmatrix} 4 & 2 \\ 2 & 1 \end{bmatrix}.
\]

\textbf{Step 1: eigenvalues.} Solve $\det(\mat{A}-\lambda\mat{I})=0$:
\[
  \det\begin{bmatrix} 4-\lambda & 2 \\ 2 & 1-\lambda \end{bmatrix}
  = (4-\lambda)(1-\lambda) - 4
  = \lambda^2 - 5\lambda + 0 = \lambda(\lambda - 5).
\]
So $\lambda_1 = 5$ and $\lambda_2 = 0$. Both real, as promised. (A zero
eigenvalue just means $\mat{A}$ flattens one direction -- it is singular.)

\textbf{Step 2: eigenvectors.} For $\lambda_1 = 5$, solve $(\mat{A}-5\mat{I})\vec{x}=\vec{0}$:
\[
  \begin{bmatrix} -1 & 2 \\ 2 & -4 \end{bmatrix}\vec{x}=\vec{0}
  \;\Rightarrow\; -x_1 + 2x_2 = 0 \;\Rightarrow\; \vec{x}=\colvec{2\\1}.
\]
For $\lambda_2 = 0$, solve $\mat{A}\vec{y}=\vec{0}$:
\[
  \begin{bmatrix} 4 & 2 \\ 2 & 1 \end{bmatrix}\vec{y}=\vec{0}
  \;\Rightarrow\; 2y_1 + y_2 = 0 \;\Rightarrow\; \vec{y}=\colvec{1\\-2}.
\]
Check orthogonality: $(2,1)\cdot(1,-2) = 2 - 2 = 0$. Perpendicular, as
Theorem~\ref{ch26:thm-orthogonal} guarantees.

\textbf{Step 3: normalize.} Both eigenvectors have length $\sqrt{5}$, so
\[
  \vec{q}_1 = \frac{1}{\sqrt5}\colvec{2\\1},\qquad
  \vec{q}_2 = \frac{1}{\sqrt5}\colvec{1\\-2}.
\]

\textbf{Step 4: assemble.}
\[
  \mat{Q} = \frac{1}{\sqrt5}\begin{bmatrix} 2 & 1 \\ 1 & -2 \end{bmatrix},
  \qquad
  \mat{\Lambda} = \begin{bmatrix} 5 & 0 \\ 0 & 0 \end{bmatrix}.
\]
Verify $\mat{Q}\T\mat{Q}=\mat{I}$:
$\frac15\begin{bmatrix}2&1\\1&-2\end{bmatrix}\begin{bmatrix}2&1\\1&-2\end{bmatrix}
=\frac15\begin{bmatrix}5&0\\0&5\end{bmatrix}=\mat{I}$. Good -- $\mat{Q}$ is
orthogonal.

\textbf{Step 5: spectral form.} As a sum of projections,
\[
  \mat{A} = 5\,\vec{q}_1\vec{q}_1\T + 0\,\vec{q}_2\vec{q}_2\T
  = 5\cdot\frac15\begin{bmatrix}2\\1\end{bmatrix}\begin{bmatrix}2&1\end{bmatrix}
  = \begin{bmatrix}4&2\\2&1\end{bmatrix}.
\]
The whole matrix is recovered from a \emph{single} rank-one projection scaled by
$5$, because the other eigenvalue is $0$. That last picture -- a symmetric matrix
as a stack of scaled projections -- is the one to remember.
\end{example}

\begin{labbox}
Diagonalize $\mat{A}=\begin{bmatrix}3&1\\1&3\end{bmatrix}$ by hand, following the
five steps above. You should find $\lambda = 4$ and $\lambda = 2$, eigenvectors
$(1,1)$ and $(1,-1)$, and
\[
  \mat{Q}=\frac{1}{\sqrt2}\begin{bmatrix}1&1\\1&-1\end{bmatrix},\quad
  \mat{\Lambda}=\begin{bmatrix}4&0\\0&2\end{bmatrix}.
\]
Then \emph{verify} by multiplying $\mat{Q}\mat{\Lambda}\mat{Q}\T$ back out and
confirming you recover $\mat{A}$. Finally, write $\mat{A}$ as
$4\,\vec{q}_1\vec{q}_1\T + 2\,\vec{q}_2\vec{q}_2\T$ and check that too. Feeling the
transpose do the work of an inverse is the whole point.
\end{labbox}

\section{The Spectral Theorem and Quadratic Forms}
\label{ch26:quadratic-forms}

Symmetric matrices are the natural language of \textbf{quadratic forms} --
expressions like $3x^2 + 4xy + 3y^2$ that are degree-two in the variables. Every
such form can be written as $q(\vec{x}) = \vec{x}\T\mat{A}\vec{x}$ for a unique
symmetric $\mat{A}$. For instance
\[
  3x^2 + 4xy + 3y^2
  = \begin{bmatrix}x&y\end{bmatrix}\begin{bmatrix}3&2\\2&3\end{bmatrix}\colvec{x\\y}.
\]
The off-diagonal $2$'s are each \emph{half} the cross-term coefficient $4$, split
evenly so the matrix comes out symmetric. (That even split is exactly why
quadratic forms force symmetry on us.)

The spectral theorem tells you what such a form ``really'' looks like. Substitute
$\mat{A}=\mat{Q}\mat{\Lambda}\mat{Q}\T$ and change to the rotated coordinates
$\vec{y}=\mat{Q}\T\vec{x}$:
\[
  q(\vec{x}) = \vec{x}\T\mat{Q}\mat{\Lambda}\mat{Q}\T\vec{x}
  = (\mat{Q}\T\vec{x})\T\mat{\Lambda}(\mat{Q}\T\vec{x})
  = \vec{y}\T\mat{\Lambda}\vec{y}
  = \lambda_1 y_1^2 + \lambda_2 y_2^2 + \cdots + \lambda_n y_n^2.
\]
In the right coordinates -- the principal axes -- the cross terms vanish entirely
and the form is a plain weighted sum of squares. The eigenvalues are the weights.

\begin{figure}[H]
\centering
\begin{tikzpicture}[scale=1.0,>=Stealth]
  \draw[->,gray!55] (-2.6,0)--(2.6,0) node[right]{$x$};
  \draw[->,gray!55] (0,-2.6)--(0,2.6) node[above]{$y$};
  % level set of 3x^2+4xy+3y^2 = 1 : tilted ellipse, rotated 45deg.
  % Semi-axis along q1=(1,1) is 1/sqrt(lambda1)=1/sqrt5~0.447 (SHORT, large eigenvalue);
  % semi-axis along q2=(1,-1) is 1/sqrt(lambda2)=1 (LONG). With rotate=45 the first
  % radius sits along (1,1), the second along (1,-1), so order is (0.447 and 1).
  \draw[targetgreen,thick,rotate=45] (0,0) ellipse (0.447 and 1);
  % principal axes: draw them to the ellipse boundary so the large eigenvalue (lambda=5)
  % visibly gives the SHORT semi-axis along q1 and lambda=1 gives the LONG one along q2.
  \draw[->,vecred,thick] (0,0)--(0.316,0.316) node[above right]{\small$\vec{q}_1\ (\lambda=5)$};
  \draw[->,structurepurple,thick] (0,0)--(0.707,-0.707) node[below right]{\small$\vec{q}_2\ (\lambda=1)$};
  \node[targetgreen] at (-1.3,1.2) {\small$q(\vec{x})=1$};
\end{tikzpicture}
\caption{The level set $3x^2+4xy+3y^2=1$ of a quadratic form is an ellipse whose
axes are the eigenvectors of $\protect\begin{bmatrix}3&2\\2&3\protect\end{bmatrix}$
(here $\lambda_1=5$ along $(1,1)$, $\lambda_2=1$ along $(1,-1)$). A \emph{larger}
eigenvalue makes the form grow faster, so the ellipse is \emph{shorter} along that
axis (semi-axis length $1/\sqrt{\lambda}$).}
\label{ch26:fig-quadform}
\end{figure}

This is the bridge to the next chapter on \emph{positive definite} matrices: the
form $q(\vec{x})$ is positive for every nonzero $\vec{x}$ exactly when all the
eigenvalues $\lambda_k$ are positive -- because in principal coordinates it is a
sum of positive multiples of squares. Symmetry plus positive eigenvalues is the
algebra of bowls that curve upward in every direction, which is the geometry of a
unique minimum, which is why optimization cares so much.

\begin{connectionbox}
\textbf{Calculus and the second derivative test.} Near a critical point, a smooth
function of several variables is approximated by a quadratic form whose matrix is
the symmetric \emph{Hessian} $H$ of second partials. The spectral theorem
diagonalizes $H$ into independent parabolas along its principal axes. The critical
point is a \emph{minimum} when all eigenvalues of $H$ are positive (a bowl), a
\emph{maximum} when all are negative (a dome), and a \emph{saddle} when the signs
disagree -- up along one principal axis, down along a perpendicular one. The
familiar single-variable second-derivative test ($f''>0$ means minimum) is just
the $1\times 1$ case, and the multivariable test is the spectral theorem in
disguise.
\end{connectionbox}

\section{Why Symmetry Is Everywhere: A Modern Application}
\label{ch26:application}

\begin{applicationbox}
\textbf{Principal Component Analysis (PCA).} Suppose you have a cloud of data --
say $n$ samples, each a point in $\R^d$ (pixels of a face, genes of a cell,
features of a customer). Center it so the mean is at the origin and stack the
samples as rows of a matrix $\mat{X}$. The \textbf{covariance matrix}
\[
  \mat{C} = \frac{1}{n-1}\,\mat{X}\T\mat{X}
\]
is symmetric (it is a Gram matrix, Section~\ref{ch26:where-from}), so the spectral
theorem applies: $\mat{C} = \mat{Q}\mat{\Lambda}\mat{Q}\T$ with real eigenvalues
$\lambda_1 \ge \lambda_2 \ge \cdots \ge 0$ and orthonormal eigenvectors
$\vec{q}_k$.

Those eigenvectors are the \textbf{principal components}, the perpendicular
directions along which the data varies. The eigenvalue $\lambda_k$ is the
\emph{variance} of the data along direction $\vec{q}_k$. The largest eigenvalue
points along the direction of greatest spread -- the long axis of the data
ellipsoid in Figure~\ref{ch26:fig-ellipse}, now drawn by real measurements rather
than a chosen matrix.

\emph{The payoff:} keep only the top few $\vec{q}_k$ (largest variances) and you
have compressed high-dimensional data onto a few perpendicular axes while losing
the least information possible. Eigenfaces, topic models, the ``elbow plot'' of
explained variance, dimensionality reduction before a neural net -- all of it is
this single computation: \emph{diagonalize a symmetric covariance matrix and read
its principal axes.} The reason PCA's directions come out conveniently
perpendicular and its variances come out conveniently real is precisely the
spectral theorem.
\end{applicationbox}

\begin{remark}[Why symmetry reflects nature]
\label{ch26:rem-not-convenience}
Symmetry is not merely a convenient hypothesis chosen for clean theorems.
Quantities such as energy, variance, distance, and mutual coupling are symmetric
by construction because they measure undirected relationships between pairs,
which is why the spectral theorem applies so widely.
\end{remark}

\section{A Gallery of Symmetric Behaviors}
\label{ch26:gallery}

A few small worked cases to anchor the range of what symmetric matrices do.

\begin{example}[Diagonal matrices: the trivial spectral case]
\label{ch26:ex-diagonal}
A diagonal matrix $\mat{D}=\diag(d_1,\dots,d_n)$ is already symmetric and already
diagonalized: $\mat{Q}=\mat{I}$, $\mat{\Lambda}=\mat{D}$. Its principal axes are
the standard coordinate axes, and it stretches axis $k$ by $d_k$. Every symmetric
matrix is just a diagonal matrix viewed in a rotated coordinate system. This is
the entire content of the spectral theorem in one sentence.
\end{example}

\begin{example}[A reflection is symmetric]
\label{ch26:ex-reflection}
The reflection across the line $y=x$ is
\[
  \mat{R} = \begin{bmatrix} 0 & 1 \\ 1 & 0 \end{bmatrix},\qquad \mat{R}\T=\mat{R}.
\]
It is symmetric, so it has perpendicular eigenvectors with real eigenvalues. Along
the mirror line $(1,1)$ nothing moves: $\mat{R}(1,1)=(1,1)$, eigenvalue $+1$.
Perpendicular to the mirror, $(1,-1)$ flips: $\mat{R}(1,-1)=(-1,1)=-(1,-1)$,
eigenvalue $-1$. So $\mat{\Lambda}=\diag(1,-1)$. The \emph{negative} eigenvalue is
the algebra of ``this axis got flipped.'' Note a \emph{rotation} (other than by
$0$ or $\pi$) is \emph{not} symmetric -- consistent with its complex eigenvalues
and its lack of any real stretching axis.
\end{example}

\begin{example}[Trace and determinant as quick checks]
\label{ch26:ex-trace-det}
Two sanity checks cost nothing. The sum of eigenvalues equals the trace, and the
product equals the determinant. For $\mat{A}=\begin{bmatrix}4&2\\2&1\end{bmatrix}$
from Example~\ref{ch26:ex-full-2x2}: $\tr\mat{A}=4+1=5=\lambda_1+\lambda_2=5+0$,
and $\det\mat{A}=4-4=0=\lambda_1\lambda_2=5\cdot 0$. Both check out. For a
$2\times 2$ symmetric matrix this lets you read off the eigenvalues almost by
inspection: they are the two numbers summing to the trace and multiplying to the
determinant.
\end{example}

\begin{trapbox}
A real matrix with real eigenvalues is \emph{not} necessarily symmetric. The shear
$\begin{bmatrix}1&1\\0&1\end{bmatrix}$ has the real (repeated) eigenvalue $1$ but
is not symmetric and is not even diagonalizable -- it has only a one-dimensional
line of eigenvectors. Symmetry is a much stronger condition than ``real
eigenvalues.'' What symmetry buys you on top of real eigenvalues is the
\emph{orthonormal basis} of eigenvectors. Do not confuse the two.
\end{trapbox}

\begin{historybox}
The Spectral Theorem was forged in the study of continuum mechanics, elasticity, and celestial perturbations. When an elastic solid is deformed under pressure, the internal forces across any microscopic plane are described by the \textbf{stress tensor} (a $3\times 3$ matrix describing the internal forces of tension, compression, and shear pushing across any cross-section of the material). Because conservation of angular momentum requires internal torques to balance, the stress tensor is strictly symmetric: $\sigma_{ij} = \sigma_{ji}$. In arbitrary coordinates, normal forces (compression and tension) are entangled with shear forces that warp right angles. In the 1750s, Leonhard Euler proved that every rotating rigid body has three perpendicular axes where rotation produces no wobble. In elasticity, the fundamental physical question was identical: does there always exist a perpendicular frame of axes where all shearing vanishes entirely, leaving only pure, independent compression or tension?

In 1829, Augustin-Louis Cauchy published the first proof of the principal axis theorem for $3\times 3$ symmetric matrices in celestial mechanics. Yet two controversies lingered for decades. First, why must the eigenvalues of a symmetric matrix be strictly real? General matrices routinely produce pairs of complex conjugate roots corresponding to rotations. Cauchy showed that if $\mat{A}\vec{z} = \lambda\vec{z}$ for a complex vector, symmetry forces $\bar{\vec{z}}\T\mat{A}\vec{z} = \lambda \norm{\vec{z}}^2$ to equal its own conjugate transpose, ensuring that every eigenvalue is real. Second, what happens when eigenvalues are repeated? General matrices with repeated eigenvalues frequently become defective (short of eigenvectors). In 1858, Karl Weierstrass delivered the complete, general proof for $n\times n$ matrices using variational methods (finding eigenvalues by maximizing or minimizing the energy ratio $\vec{x}\T\mat{A}\vec{x}/\norm{\vec{x}}^2$, the Rayleigh quotient, rather than factoring polynomials). He proved that a real symmetric matrix can never be defective: its geometric multiplicity always matches its algebraic multiplicity (the dimension of each eigenspace matches how many times the eigenvalue repeats as a polynomial root), guaranteeing an orthonormal basis of eigenvectors under all circumstances.

What the Spectral Theorem delivers is the cleanest factorization in linear algebra: $\mat{A} = \mat{Q}\mat{\Lambda}\mat{Q}\T$, expressing any symmetric operator as a pure stretch along perpendicular directions with an inverse that costs nothing ($\mat{Q}\inv = \mat{Q}\T$). Its honest limitation is that it applies only to symmetric matrices; real-world data matrices and directed networks are rarely symmetric, and its eigenvalues may still be negative, creating saddle points (surfaces that curve up along one axis and down along another) instead of energy bowls. In Chapter~27, we restrict our focus to symmetric matrices with strictly positive eigenvalues, and in Chapter~29, we liberate the Spectral Theorem from the symmetry restriction entirely, applying it to $\mat{A}\T\mat{A}$ to reach the singular value decomposition for every rectangular matrix. In infinite dimensions (in quantum mechanics, where physical observables are self-adjoint differential operators), the spectrum can feature continuous unbroken bands rather than discrete eigenvalues, leading to the functional-analytic spectral theorem developed by David Hilbert and John von Neumann; the statistical behavior of eigenvalues in large random symmetric matrices remains an active frontier in mathematical physics.
\end{historybox}

\section{Summary}
\label{ch26:summary}

A symmetric matrix $\mat{A}=\mat{A}\T$ is, geometrically, nothing but a stretch
along a set of perpendicular axes. The dot-product identity
$\inner{\mat{A}\vec{x}}{\vec{y}}=\inner{\vec{x}}{\mat{A}\vec{y}}$ is the engine
that makes everything work:

\begin{itemize}
  \item \textbf{Real eigenvalues.} No twisting; only stretching
        (Theorem~\ref{ch26:thm-real}).
  \item \textbf{Orthogonal eigenvectors.} Distinct eigenvalues force perpendicular
        directions, and repeated ones can always be chosen perpendicular
        (Theorem~\ref{ch26:thm-orthogonal}).
  \item \textbf{Spectral theorem.} $\mat{A}=\mat{Q}\mat{\Lambda}\mat{Q}\T$ with
        $\mat{Q}$ orthogonal and $\mat{\Lambda}$ real diagonal
        (Theorem~\ref{ch26:thm-spectral}); equivalently
        $\mat{A}=\sum_k \lambda_k\,\vec{q}_k\vec{q}_k\T$, a weighted sum of
        perpendicular projections.
  \item \textbf{Geometry.} Rotate the principal axes onto the coordinate axes,
        stretch by the eigenvalues, rotate back -- with the inverse rotation
        delivered for free by a transpose.
  \item \textbf{Cost.} The orthonormal basis provides a major computational advantage:
        $\mat{Q}\inv=\mat{Q}\T$, so applying $\mat{A}$ or any matrix function of it
        ($\mat{A}^k$, $\sqrt{\mat{A}}$, $\mat{A}\inv$) reduces to ``transpose, scale
        the eigenvalues, transpose back'' -- with no general matrix inversion required,
        and orthogonal steps that maintain numerical stability.
  \item \textbf{Why it matters.} Covariances, Hessians, Gram matrices, distances,
        and undirected networks are all symmetric, so PCA, the second-derivative
        test, least squares, and the SVD all stand on this theorem.
\end{itemize}

In Chapter~27, we specialize one notch further: what happens when every
eigenvalue is strictly positive? The ellipse becomes a true energy ellipsoid,
the quadratic form $q(\vec{x}) = \vec{x}\T\mat{A}\vec{x}$ curves upward in
every direction like a bowl with a unique minimum at the origin, and the matrix
earns the title \textbf{positive definite}. That unlocks the Cholesky
factorization $\mat{A} = \mat{R}\T\mat{R}$, generalized dot products, and the
computational engine of convex optimization.

\section{Exercises}
\label{ch26:exercises}

\begin{exercise}[Warmup: spot the symmetric ones]
\label{ch26:ex-spot}
Which of these are symmetric? For each symmetric one, state without computing what
its principal axes will be (perpendicular, real eigenvalues) and whether you can
predict any eigenvalue by inspection.
\[
  \begin{bmatrix}1&3\\3&1\end{bmatrix},\quad
  \begin{bmatrix}1&3\\-3&1\end{bmatrix},\quad
  \begin{bmatrix}2&0&0\\0&5&0\\0&0&5\end{bmatrix},\quad
  \begin{bmatrix}0&1\\1&0\end{bmatrix}.
\]
\end{exercise}

\begin{exercise}[Full decomposition by hand]
\label{ch26:ex-byhand}
Find the eigenvalues, orthonormal eigenvectors, $\mat{Q}$, and $\mat{\Lambda}$ for
$\mat{A}=\begin{bmatrix}5&-2\\-2&5\end{bmatrix}$. Verify $\mat{Q}\T\mat{Q}=\mat{I}$
and that $\mat{Q}\mat{\Lambda}\mat{Q}\T$ reproduces $\mat{A}$. Then write $\mat{A}$
as a sum $\lambda_1\vec{q}_1\vec{q}_1\T + \lambda_2\vec{q}_2\vec{q}_2\T$.
\end{exercise}

\begin{exercise}[Gram matrices are always symmetric]
\label{ch26:ex-gram}
Let $\mat{M}$ be any (possibly rectangular) real matrix. Prove
$\mat{M}\T\mat{M}$ is symmetric using only transpose rules. Then compute it for
$\mat{M}=\begin{bmatrix}1&2\\0&1\\1&0\end{bmatrix}$ and confirm the result is
symmetric.
\end{exercise}

\begin{exercise}[The orthogonality proof, in your own hand]
\label{ch26:ex-ortho-proof}
Let $\mat{A}$ be symmetric with $\mat{A}\vec{x}=\lambda\vec{x}$ and
$\mat{A}\vec{y}=\mu\vec{y}$. Starting from the single quantity
$\inner{\mat{A}\vec{x}}{\vec{y}}$, derive $(\lambda-\mu)\inner{\vec{x}}{\vec{y}}=0$
and explain in one sentence why $\lambda\neq\mu$ forces orthogonality. Where
exactly did you use $\mat{A}\T=\mat{A}$?
\end{exercise}

\begin{exercise}[Quadratic form to matrix and back]
\label{ch26:ex-quadform}
Write $q(x,y)=2x^2 - 6xy + 2y^2$ as $\vec{x}\T\mat{A}\vec{x}$ with $\mat{A}$
symmetric. Diagonalize $\mat{A}$, find its principal axes, and describe the level
set $q(\vec{x})=1$ (ellipse? hyperbola?). What about the signs of the eigenvalues
makes it one or the other?
\end{exercise}

\begin{exercise}[Trace and determinant shortcut]
\label{ch26:ex-trace}
A symmetric $2\times 2$ matrix has trace $6$ and determinant $5$. Without finding
the matrix, give its eigenvalues. Are its eigenvectors guaranteed perpendicular?
Why?
\end{exercise}

\begin{exercise}[Reflections and signs]
\label{ch26:ex-reflection-ex}
The matrix $\mat{R}=\frac{1}{5}\begin{bmatrix}3&4\\4&-3\end{bmatrix}$ is symmetric.
Show $\mat{R}\T\mat{R}=\mat{I}$ (so $\mat{R}$ is also orthogonal). Find its
eigenvalues -- you should get $+1$ and $-1$ -- and interpret $\mat{R}$ geometrically
as a reflection. Along which line does the $+1$ eigenvector point?
\end{exercise}

\begin{exercise}[Functions of a symmetric matrix]
\label{ch26:ex-function}
Using $\mat{A}=\mat{Q}\mat{\Lambda}\mat{Q}\T$, explain why
$\mat{A}^k=\mat{Q}\mat{\Lambda}^k\mat{Q}\T$, and why one may sensibly define
$\sqrt{\mat{A}}=\mat{Q}\sqrt{\mat{\Lambda}}\,\mat{Q}\T$ (taking square roots of the
eigenvalues) when all $\lambda_k\ge 0$. Compute $\sqrt{\mat{A}}$ for
$\mat{A}=\begin{bmatrix}2&1\\1&2\end{bmatrix}$ (eigenvalues $3$ and $1$). What goes
wrong if some eigenvalue is negative?
\end{exercise}

\begin{exercise}[The invariance step, made concrete]
\label{ch26:ex-invariance}
The proof sketch hinged on: if $\vec{q}_1$ is an eigenvector of symmetric
$\mat{A}$ and $\vec{w}\perp\vec{q}_1$, then $\mat{A}\vec{w}\perp\vec{q}_1$ too.
Verify this concretely for $\mat{A}=\begin{bmatrix}2&1\\1&2\end{bmatrix}$,
$\vec{q}_1=(1,1)$, and $\vec{w}=(1,-1)$: check $\vec{w}\perp\vec{q}_1$, compute
$\mat{A}\vec{w}$, and confirm it is still perpendicular to $\vec{q}_1$.
\end{exercise}

\begin{exercise}[Challenge: a symmetric matrix from a rotation]
\label{ch26:ex-challenge}
Let $\mat{B}=\begin{bmatrix}\cos\theta&-\sin\theta\\\sin\theta&\cos\theta\end{bmatrix}$
be a rotation. It is \emph{not} symmetric (for generic $\theta$). But
$\mat{B}\T\mat{B}$ is symmetric -- compute it and identify the result. What are its
eigenvalues, and what does that say about how a rotation stretches lengths?
Connect your answer to the fact that rotations are rigid motions.
\end{exercise}

\begin{exercise}[Challenge: spectral theorem fails without symmetry]
\label{ch26:ex-fail}
Give a $2\times 2$ real matrix that has two real eigenvalues but whose
eigenvectors are \emph{not} perpendicular, demonstrating that orthogonality is a
genuine consequence of symmetry and not of real eigenvalues alone. (Hint: try an
upper-triangular matrix with distinct diagonal entries.) Then give a $2\times 2$
real matrix with \emph{no} real eigenvalues at all, and explain why it cannot be
symmetric.
\end{exercise}
